\section{Build the Python environment}
\label{sec:python}

\subsection{Background}

The project depends on nine Python \term{packages} (libraries written by others), among them
NumPy for arithmetic on arrays, OpenCV for image processing and TensorFlow for the neural
network; the file \file{requirements.txt} lists them. They are installed into a \term{virtual
environment}: a private copy of Python that lives in a folder named \file{.venv} inside the
project and is used by this project alone. Deleting that folder removes the environment
completely, and no other program on the laptop is affected by what is installed in it.

A PowerShell window uses the environment only after the environment has been \term{activated}
in that window. Activation lasts until the window is closed, and it is shown by the word
\texttt{(.venv)} at the start of the prompt.

\subsection{Create and activate the environment}

\step{Start in the project folder}
Use a PowerShell window whose prompt reads
\plainprompt
If the prompt shows another folder, run:
\begin{ps}
cd C:\ee105\camera-to-latex
\end{ps}

\step{Create the environment}
\begin{ps}
py -3.12 -m venv .venv
\end{ps}
The command works for about half a minute and prints nothing.
\checkpoint{A new prompt has appeared without any red text, and the command
\texttt{Get-ChildItem} now lists a folder named \file{.venv}.}

\step{Allow PowerShell to run the activation script}\label{st:policy}
A new Windows installation refuses to run script files, and activation is a script. The next
command permits scripts that are stored on your own computer, for your user account only.
\begin{ps}
Set-ExecutionPolicy RemoteSigned -Scope CurrentUser
\end{ps}
PowerShell explains the change and asks for confirmation. Type \code|Y| and press \key{Enter}.
\checkpoint{A new prompt has appeared. The command is needed once per computer; it does no harm
to run it again.}

\step{Activate the environment}\label{st:activate}
\begin{ps}
.\.venv\Scripts\Activate.ps1
\end{ps}
The line begins with a full stop and a backslash, which mean ``in the current folder''.
\begin{out}
(.venv) PS C:\ee105\camera-to-latex>
\end{out}
\checkpoint{The prompt now begins with \texttt{(.venv)}.}

\begin{iffail}[If red text says that running scripts is disabled]
Step~\ref{st:policy} did not take effect. Run it again, answer \code|Y|, and repeat this step.
If the laptop is managed by an organization that forbids the change,
Section~\ref{sec:trouble} gives an alternative.
\end{iffail}

\subsection{Install the packages}

\step{Update the installer program}
The program that installs Python packages is called \cmd{pip}. Bring it up to date first:
\begin{ps}
python -m pip install --upgrade pip
\end{ps}
\checkpoint{The output ends with \emph{Successfully installed pip} and a version number, or
says that the requirement is already satisfied.}

\step{Install the project's packages}\label{st:pipinstall}
\begin{ps}
pip install -r requirements.txt
\end{ps}
This is the longest installation of the guide: between 10 and 30 minutes, depending on the
internet connection and on the antivirus program, which inspects every file. The output scrolls
through lines that begin with \emph{Collecting} and \emph{Downloading}; the largest download,
TensorFlow, is about 350\,MB. Near the end, the window shows a line that begins with
\emph{Installing collected packages} and then appears to stand still, often for 5 to 15 minutes.
\textbf{This is normal}: pip is unpacking many thousands of small files and reports nothing
while it does so. Leave the window alone.
\begin{out}[What you should see at the end (the list is abridged here)]
Installing collected packages: pyserial, namex, libclang, flatbuffers, ..., keras, tensorflow
Successfully installed ... keras-3.15.1 ... tensorflow-2.21.0 ...
(.venv) PS C:\ee105\camera-to-latex>
\end{out}
\checkpoint{A line beginning with \emph{Successfully installed} names \emph{tensorflow}, and a
new prompt has appeared. A notice that a newer version of pip exists can be ignored.}

\begin{iffail}[If the window seems frozen]
Look at the title bar. If it contains the word \emph{Select}, a mouse click has paused the
output: press \key{Esc}. Otherwise wait. To convince yourself that work continues, open Task
Manager with \key{Ctrl}+\key{Shift}+\key{Esc} and look for \emph{Python} in the list of
processes with a changing CPU or disk figure. Interrupt the command (\key{Ctrl}+\key{C}) only if
nothing at all has changed for 30 minutes, and then simply run the same command again; pip
continues with what is missing.
\end{iffail}

\step{Fix the version of OpenCV}\label{st:opencv}
\begin{ps}
pip install "opencv-python>=4.8,<5"
\end{ps}
\rationale{Version~5 of the image-processing library OpenCV removed a function
(\cmd{HOGDescriptor}) that the project's fallback classifier uses. The command makes sure that a
version~4 release is installed; the quotation marks are required because PowerShell would
otherwise interpret the signs \texttt{<} and \texttt{>} itself.}
\checkpoint{The output ends either with \emph{Successfully installed opencv-python-4} followed
by more digits, or with \emph{Requirement already satisfied}.}

\step{Verify the two most important packages}\label{st:verifypy}
\begin{ps}
python -c "import cv2, tensorflow as tf; print(cv2.__version__, tf.__version__)"
\end{ps}
The command takes 10 to 30 seconds the first time. TensorFlow prints a few informational lines
first; they mention \emph{oneDNN} and are not errors.
\begin{out}[What you should see (last line)]
4.12.0 2.21.0
\end{out}
\checkpoint{The last line shows two version numbers; the first begins with \texttt{4.} and the
second with \texttt{2.}}

\subsection{The opening sequence for every later session}
\label{sec:opening}

The environment has to be activated again in every new PowerShell window. From now on, whenever
this guide or its companion asks for a new window, open one as in Step~\ref{st:openps} and run
the following two commands. Together they are called the \term{opening sequence}.

\begin{ps}[Opening sequence, command 1 of 2]
cd C:\ee105\camera-to-latex
\end{ps}
\begin{ps}[Opening sequence, command 2 of 2]
.\.venv\Scripts\Activate.ps1
\end{ps}
\checkpoint{The prompt reads as follows. Every command of the project is run from this folder
with this prefix showing.}
\venvprompt

\begin{note}[Optional: silence TensorFlow's informational lines]
The command \code|$env:TF_CPP_MIN_LOG_LEVEL = "2"| hides most of the lines that TensorFlow
prints when it starts. It affects only the window in which it is typed, and nothing depends on
it.
\end{note}
